\documentclass[11pt,a4paper]{article}

% ------------------------------------------------------------------
% Encodage & langue
% ------------------------------------------------------------------
\usepackage[utf8]{inputenc}
\usepackage[T1]{fontenc}
\usepackage[french]{babel}
\usepackage{lmodern}

% ------------------------------------------------------------------
% Mise en page sobre
% ------------------------------------------------------------------
\usepackage[a4paper,margin=2.5cm]{geometry}
\usepackage{microtype}
\usepackage{setspace}
\onehalfspacing

% ------------------------------------------------------------------
% Couleurs minimales (niveaux de gris uniquement)
% ------------------------------------------------------------------
\usepackage{xcolor}
\definecolor{gris}{gray}{0.45}
\definecolor{grisclair}{gray}{0.92}
\definecolor{grisfonce}{gray}{0.20}

% ------------------------------------------------------------------
% Tableaux, listings, liens
% ------------------------------------------------------------------
\usepackage{booktabs}
\usepackage{tabularx}
\usepackage{array}
\usepackage{longtable}
\usepackage{enumitem}
\setlist{itemsep=2pt,topsep=3pt}

\usepackage{listings}
\lstset{
  basicstyle=\ttfamily\footnotesize,
  backgroundcolor=\color{grisclair},
  frame=single,
  rulecolor=\color{gris},
  framerule=0.3pt,
  breaklines=true,
  showstringspaces=false,
  columns=fullflexible,
  keepspaces=true,
  xleftmargin=4pt,
  xrightmargin=4pt,
  aboveskip=8pt,
  belowskip=8pt,
}

\usepackage[hidelinks]{hyperref}
\hypersetup{
  colorlinks=true,
  linkcolor=grisfonce,
  urlcolor=grisfonce,
  citecolor=grisfonce,
}

% ------------------------------------------------------------------
% En-têtes / pieds de page
% ------------------------------------------------------------------
\usepackage{fancyhdr}
\pagestyle{fancy}
\fancyhf{}
\renewcommand{\headrulewidth}{0.3pt}
\renewcommand{\footrulewidth}{0pt}
\lhead{\small\color{gris} Pipeline de données en temps réel}
\rhead{\small\color{gris} FinHub Pipeline}
\cfoot{\small\color{gris}\thepage}

% Titres de section discrets
\usepackage{titlesec}
\titleformat{\section}{\Large\bfseries\color{grisfonce}}{\thesection}{0.6em}{}
\titleformat{\subsection}{\large\bfseries\color{grisfonce}}{\thesubsection}{0.6em}{}
\titleformat{\subsubsection}{\normalsize\bfseries\color{grisfonce}}{\thesubsubsection}{0.6em}{}

% Commande pour les chemins / identifiants
\newcommand{\code}[1]{\texttt{\small #1}}

% ==================================================================
\begin{document}
% ==================================================================

% ------------------------------------------------------------------
% Page de titre sobre
% ------------------------------------------------------------------
\begin{titlepage}
\centering
\vspace*{3cm}

{\color{gris}\rule{\linewidth}{0.4pt}}\\[0.6cm]
{\Huge\bfseries\color{grisfonce} Pipeline de données\\[0.2cm] en temps réel}\\[0.5cm]
{\large\color{gris} Données du marché boursier : de l'ingestion à la modélisation analytique}\\[0.4cm]
{\color{gris}\rule{\linewidth}{0.4pt}}\\[1.5cm]

{\large\textbf{FinHub Streaming Data Pipeline}}\\[0.3cm]
{\normalsize\color{gris} Finnhub $\rightarrow$ Kafka $\rightarrow$ Data Lake $\rightarrow$ dbt}\\[3cm]

\begin{tabular}{rl}
\textbf{Auteur} & Salim Diallo \\[4pt]
\textbf{Dépôt} & \code{PipelineFinhubStreamingData} \\[4pt]
\textbf{Technologies} & Python, Kafka, MinIO (S3), Airflow, dbt, Docker \\[4pt]
\textbf{Date} & Juin 2026 \\
\end{tabular}

\vfill
{\small\color{gris} Rapport technique de conception et de réalisation}
\end{titlepage}

% ------------------------------------------------------------------
\tableofcontents
\newpage

% ==================================================================
\section{Introduction}
% ==================================================================

\subsection{Contexte et objectif}

Ce projet consiste à concevoir une \textbf{pipeline de données en temps réel}
exploitant les données du marché boursier. L'objectif est de récupérer des
données financières en direct, de les traiter par couches successives, puis de
les mettre à disposition pour des analyses fiables et reproductibles.

La chaîne de traitement va du \emph{streaming} temps réel jusqu'à la
modélisation analytique, en respectant les bonnes pratiques d'ingénierie de
données : séparation des responsabilités, couches de qualité croissante
(\emph{medallion architecture}), orchestration et conteneurisation.

\subsection{Vue d'ensemble du flux}

\begin{center}
\fbox{\parbox{0.9\linewidth}{\centering\ttfamily\small
Finnhub API (WebSocket) \\
$\downarrow$ \\
Kafka (topics par type d'événement) \\
$\downarrow$ \\
S3 / MinIO -- couche Bronze (brut) \\
$\downarrow$ \\
Couche Silver (nettoyé, typé) \\
$\downarrow$ \\
Couche Gold (métriques business) \\
$\downarrow$ \\
dbt / DuckDB (modèles analytiques) \\
}}
\end{center}

\subsection{Synthèse des étapes}

\begin{center}
\renewcommand{\arraystretch}{1.3}
\begin{tabularx}{\linewidth}{@{}clX@{}}
\toprule
\textbf{\#} & \textbf{Étape} & \textbf{Livrable principal} \\
\midrule
1 & Ingestion \& Streaming & Client Finnhub WebSocket $\rightarrow$ producteur Kafka \\
2 & Stockage (Data Lake)   & Consumers $\rightarrow$ Bronze ; transformations Silver et Gold \\
3 & Orchestration          & DAG Apache Airflow (traitement batch) \\
4 & Modélisation           & Modèles dbt (fait, dimension, agrégats) sur DuckDB \\
5 & Infrastructure         & Conteneurisation Docker et orchestration des stacks \\
\bottomrule
\end{tabularx}
\end{center}

% ==================================================================
\section{Architecture générale}
% ==================================================================

\subsection{Principes retenus}

\begin{itemize}
  \item \textbf{Découplage} : la source (Finnhub) et les consommateurs sont
        découplés par Kafka, qui sert de tampon et de point de distribution.
  \item \textbf{Architecture medallion} : trois couches de qualité croissante
        (Bronze $\rightarrow$ Silver $\rightarrow$ Gold) dans le data lake.
  \item \textbf{Streaming vs batch} : l'ingestion est un processus continu ;
        les transformations et la modélisation sont orchestrées en batch.
  \item \textbf{Configuration centralisée} : tous les paramètres sensibles
        passent par un fichier \code{.env} (jamais versionné).
\end{itemize}

\subsection{Arborescence du projet}

\begin{lstlisting}
finHubPipeline/
|-- ingestion/          # Finnhub -> Kafka (producteur)
|   |-- finnhub_client.py
|   |-- producer.py
|   |-- config.py
|   |-- run_ingestion.py
|   `-- schemas/        # contrats JSON (trades, quotes, candles)
|-- kafka/              # stack Kafka (KRaft) + creation des topics
|-- consumers/          # Kafka -> S3 Bronze
|   |-- consumer_trades.py / quotes / candles
|   `-- utils.py
|-- transformations/
|   |-- silver/         # nettoyage -> Parquet typé
|   `-- gold/           # metriques (VWAP, volatilite, ...)
|-- airflow/            # DAG d'orchestration + stack Airflow
|-- warehouse/dbt/      # modeles dbt (DuckDB)
|-- infra/              # Dockerfile + MinIO + services conteneurises
|-- Makefile            # orchestration de toute l'infra
`-- pyproject.toml      # dependances (uv)
\end{lstlisting}

\subsection{Pile technologique}

\begin{center}
\renewcommand{\arraystretch}{1.25}
\begin{tabularx}{\linewidth}{@{}lX@{}}
\toprule
\textbf{Composant} & \textbf{Rôle} \\
\midrule
Python ($\geq$ 3.14) + uv & Langage et gestionnaire de dépendances \\
\code{finnhub-python}, \code{websocket-client} & Connexion à la source de données \\
\code{kafka-python} & Production et consommation Kafka \\
Apache Kafka (KRaft) & Bus de streaming (sans ZooKeeper) \\
MinIO + \code{boto3} & Data lake compatible S3 \\
\code{pandas}, \code{pyarrow} & Transformations Silver / Gold (Parquet) \\
Apache Airflow & Orchestration du flux batch \\
dbt + DuckDB & Modélisation analytique \\
Docker / Docker Compose & Conteneurisation et déploiement \\
\bottomrule
\end{tabularx}
\end{center}

% ==================================================================
\section{Étape 1 — Ingestion \& Streaming}
% ==================================================================

\subsection{Source de données : Finnhub}

Les données proviennent de l'API \textbf{Finnhub}. Trois types d'événements
sont modélisés dans le projet : \emph{trades} (transactions), \emph{quotes}
(bid/ask) et \emph{candles} (OHLC). Sur le plan gratuit, seul le flux
\emph{trade} est diffusé en temps réel via WebSocket ; les quotes et candles
temps réel relèvent d'un abonnement payant. La pipeline alimente donc
concrètement le flux des trades, les autres topics restant prêts à l'emploi.

\subsection{Client Finnhub}

Le module \code{ingestion/finnhub\_client.py} expose deux modes d'accès :
\begin{itemize}
  \item \textbf{WebSocket} (\code{start\_stream}) : abonnement aux symboles et
        normalisation de chaque trade en dictionnaire
        (\code{symbol}, \code{price}, \code{volume}, \code{timestamp},
        \code{type}). La reconnexion automatique est gérée.
  \item \textbf{REST} (\code{get\_candles}) : récupération de chandeliers OHLC
        (endpoint premium), conservé pour un usage futur.
\end{itemize}

\subsection{Production vers Kafka}

Le producteur \code{ingestion/producer.py} publie chaque événement avec :
\begin{itemize}
  \item \textbf{clé} = \code{symbol} : garantit l'ordre des messages d'un même
        symbole sur une partition ;
  \item \textbf{valeur} = JSON de l'événement ;
  \item \textbf{topic} déterminé par le type (\code{stock.trades},
        \code{stock.quotes}, \code{stock.candles}).
\end{itemize}

Les réglages de fiabilité retenus sont \code{acks=all} (durabilité),
\code{retries=3} et \code{linger\_ms=50} (micro-batching pour le débit).

\subsection{Stack Kafka}

Kafka est déployé en \textbf{mode KRaft} (sans ZooKeeper, déprécié), via l'image
officielle \code{apache/kafka}. Une interface \emph{Kafka UI} permet de
visualiser topics, partitions et messages. Le script
\code{kafka/create\_topics.sh} crée les trois topics (3 partitions chacun) de
manière idempotente.

\subsection{Validation}

L'ingestion a été validée de bout en bout : connexion WebSocket établie,
trades temps réel reçus puis produits dans \code{stock.trades}. La vérification
côté broker a confirmé que tous les messages d'un même symbole se retrouvent sur
la même partition (clé = symbole), validant la garantie d'ordre.

% ==================================================================
\section{Étape 2 — Stockage (Data Lake)}
% ==================================================================

\subsection{Couche Bronze — données brutes}

Les consumers (\code{consumers/consumer\_*.py}) consomment les topics Kafka et
écrivent les messages \textbf{bruts} dans MinIO, sans transformation. Le format
retenu est \textbf{JSON Lines}, partitionné par date :

\begin{lstlisting}
s3://datalake/bronze/<dataset>/date=YYYY-MM-DD/<horodatage>.jsonl
\end{lstlisting}

Le consumer fonctionne par \textbf{micro-batch} (flush dès \code{BATCH\_SIZE}
messages ou après \code{BATCH\_TIMEOUT\_S} secondes). Les offsets ne sont
validés qu'\emph{après} écriture réussie dans S3, garantissant une livraison
\textbf{at-least-once} (pas de perte ; doublons éventuels dédupliqués en Silver).

\subsection{Couche Silver — données nettoyées}

Le module \code{transformations/silver/} lit le Bronze et produit du
\textbf{Parquet typé} :
\begin{itemize}
  \item typage des colonnes (\code{symbol} chaîne, montants flottants,
        \code{timestamp} entier) ;
  \item normalisation des horodatages (epoch ms $\rightarrow$ \code{event\_time}
        en datetime UTC) ;
  \item suppression des lignes aux champs critiques manquants ;
  \item déduplication.
\end{itemize}

\subsection{Couche Gold — métriques business}

Le module \code{transformations/gold/} calcule les agrégats analytiques :
\begin{itemize}
  \item \textbf{\code{agg\_stock\_metrics}} (par symbole / jour) :
        prix moyen, VWAP, volume total, volatilité, rendement journalier ;
  \item \textbf{\code{agg\_5min}} : agrégats par fenêtre de 5 minutes avec
        variation en pourcentage ;
  \item \textbf{\code{features}} : moyenne mobile, rendement, volatilité
        glissante (orienté apprentissage automatique).
\end{itemize}

\subsection{Format Bronze : choix de conception}

Le format JSON Lines a été préféré au Parquet pour la couche Bronze afin de
rester fidèle à la donnée brute (principe \emph{« Bronze = raw »}). La
conversion en Parquet et le typage interviennent en Silver, conformément à
l'architecture medallion.

\subsection{Validation}

La chaîne complète a été vérifiée sur données réelles : le flux temps réel
alimente le Bronze, puis \code{clean\_trades} produit le Silver (par exemple
343 lignes brutes ramenées à 230 lignes après déduplication, sans valeur nulle
et avec un typage correct), et la couche Gold génère les tables de métriques.

% ==================================================================
\section{Étape 3 — Orchestration (Airflow)}
% ==================================================================

\subsection{Modèle d'orchestration}

Le \emph{streaming} Finnhub $\rightarrow$ Kafka est un processus continu, qui
tourne \textbf{hors Airflow}. Airflow orchestre le traitement \textbf{batch
périodique}, ce qui correspond à son modèle d'exécution. Le DAG
\code{stock\_pipeline} s'exécute selon une planification horaire
(\code{@hourly}).

\subsection{Tâches du DAG}

\begin{center}
\renewcommand{\arraystretch}{1.25}
\begin{tabularx}{\linewidth}{@{}lX@{}}
\toprule
\textbf{Tâche} & \textbf{Action} \\
\midrule
\code{consume\_to\_s3}   & Draine les topics Kafka vers la couche Bronze (mode batch) \\
\code{transform\_silver} & Bronze $\rightarrow$ Silver (nettoyage, Parquet) \\
\code{transform\_gold}   & Silver $\rightarrow$ Gold (métriques) \\
\code{dbt\_build}        & Exécute \code{dbt run} puis \code{dbt test} \\
\bottomrule
\end{tabularx}
\end{center}

Les dépendances sont strictement séquentielles :
\[
\texttt{consume\_to\_s3} \rightarrow \texttt{transform\_silver}
\rightarrow \texttt{transform\_gold} \rightarrow \texttt{dbt\_build}
\]

\subsection{Adaptation du consumer}

Pour le mode batch, un paramètre \code{run\_once} a été ajouté au consumer : il
draine les messages disponibles puis s'arrête, au lieu de boucler
indéfiniment. Un \code{poll()} initial force l'assignation des partitions afin
d'éviter les erreurs de rééquilibrage en exécution courte.

\subsection{Déploiement}

Airflow est déployé avec l'exécuteur \code{LocalExecutor} et une base de
métadonnées PostgreSQL. Le code du projet est monté dans le conteneur et les
dépendances installées dans l'image. L'interface est accessible sur le port
local~8082.

\subsection{Validation}

Le DAG a été chargé sans erreur d'import et chaque tâche exécutée individuellement
contre l'infrastructure réelle : drainage de plus d'un millier de messages vers
le Bronze, transformation Silver et Gold, puis exécution dbt — l'ensemble de la
chaîne s'est déroulé avec succès.

% ==================================================================
\section{Étape 4 — Modélisation analytique (dbt)}
% ==================================================================

\subsection{Choix de l'entrepôt : DuckDB}

L'architecture cible mentionne Snowflake, un entrepôt cloud payant. Pour
permettre une exécution réelle, gratuite et locale, le projet utilise
\textbf{DuckDB} via l'adaptateur \code{dbt-duckdb}. DuckDB lit directement les
fichiers Parquet de la couche Silver depuis MinIO grâce à l'extension
\code{httpfs} (API S3). Les modèles SQL restent identiques à ceux qu'on
écrirait pour Snowflake ; seul l'adaptateur change.

\subsection{Modèles}

\begin{center}
\renewcommand{\arraystretch}{1.25}
\begin{tabularx}{\linewidth}{@{}llX@{}}
\toprule
\textbf{Modèle} & \textbf{Type} & \textbf{Contenu} \\
\midrule
\code{stg\_trades}        & vue   & Staging des trades Silver (typé) \\
\code{fact\_trades}       & table & Fait : un enregistrement par trade \\
\code{dim\_symbol}        & table & Dimension symbole (un par ticker) \\
\code{agg\_stock\_metrics}& table & Métriques par symbole / jour \\
\bottomrule
\end{tabularx}
\end{center}

La table \code{agg\_stock\_metrics} reproduit en SQL les métriques de la couche
Gold (VWAP, volatilité, rendement journalier). Les colonnes
\code{sector} et \code{company\_name} de \code{dim\_symbol} sont présentes par
conformité au schéma, mais nulles : ces métadonnées ne sont pas fournies par le
plan gratuit Finnhub.

\subsection{Tests}

Des tests dbt (\code{not\_null}, \code{unique}) valident l'intégrité des
modèles : non-nullité des clés et des mesures critiques, unicité du symbole dans
la dimension.

\subsection{Isolation de l'environnement}

\code{dbt-core} ne prenant pas encore en charge Python~3.14 (version du projet),
dbt est isolé dans un environnement virtuel Python~3.12 dédié pour l'usage
local. Pour l'orchestration, \code{dbt-duckdb} est installé directement dans
l'image Airflow et invoqué via \code{python -m dbt.cli.main}.

\subsection{Validation}

L'exécution a été réalisée en conditions réelles : \code{dbt run} matérialise
les quatre modèles et \code{dbt test} exécute onze tests, tous au vert, sur les
données issues de la couche Silver. La même exécution a été reproduite avec
succès depuis Airflow (tâche \code{dbt\_build}).

% ==================================================================
\section{Étape 5 — Infrastructure (Docker)}
% ==================================================================

\subsection{Image du projet}

Le fichier \code{infra/Dockerfile} construit l'image \code{finhub-pipeline}
(Python~3.14 + uv) embarquant les modules \code{ingestion}, \code{consumers} et
\code{transformations}. Elle sert à la fois au producteur et aux consumers.

\subsection{Organisation des stacks}

L'infrastructure est répartie en trois stacks Docker Compose, orchestrées par un
\code{Makefile} à la racine :

\begin{center}
\renewcommand{\arraystretch}{1.25}
\begin{tabularx}{\linewidth}{@{}lX@{}}
\toprule
\textbf{Stack} & \textbf{Services} \\
\midrule
\code{kafka/}   & Kafka (KRaft) + Kafka UI \\
\code{infra/}   & MinIO + producteur + consumer conteneurisés \\
\code{airflow/} & Airflow (webserver, scheduler) + PostgreSQL \\
\bottomrule
\end{tabularx}
\end{center}

Le \code{Makefile} fournit des cibles simples : \code{make up} (démarre tout),
\code{make down}, \code{make ps}, \code{make topics}, \code{make logs}.

\subsection{Réseaux}

Les services communiquent par les réseaux Docker partagés en utilisant les
listeners internes (\code{kafka:29092}, \code{minio:9000}), ce qui est plus
fiable que de passer par la passerelle de l'hôte.

\subsection{Validation et limite rencontrée}

L'image a été construite et les conteneurs (producteur, consumer) démarrent et
se connectent correctement à Kafka et MinIO via les réseaux internes. La seule
limite observée est liée à l'environnement de test : la saturation du disque
hôte a déclenché un refus d'écriture de MinIO (\code{XMinioStorageFull}). Il
s'agit d'une contrainte matérielle de l'environnement et non d'un défaut de la
pipeline.

% ==================================================================
\section{Démarche d'ingénierie}
% ==================================================================

\subsection{Gestion de version}

Le projet suit un modèle de branches \code{main} / \code{dev} / \code{feature/*}
et la convention \emph{Conventional Commits} (\code{feat}, \code{fix},
\code{chore}, \code{docs}\dots). Chaque étape a été développée sur sa branche,
intégrée dans \code{dev}, puis fusionnée dans \code{main}.

\subsection{Sécurité}

Le fichier \code{.env} contenant la clé API a été retiré du suivi Git et un
\code{.env.example} sans secret le documente. Un \code{.gitignore} exclut les
artefacts de runtime (environnements virtuels, journaux, bases locales,
données).

\subsection{Validation systématique}

Chaque étape a été vérifiée par une exécution réelle (et non seulement par
analyse du code) : flux de données effectif, objets écrits dans le data lake,
tâches Airflow exécutées, tests dbt au vert. Les difficultés rencontrées
(permissions Docker, résolution réseau Kafka, rééquilibrage de consumer,
compatibilité de versions) ont été diagnostiquées jusqu'à leur cause racine et
documentées.

% ==================================================================
\section{Limites et perspectives}
% ==================================================================

\subsection{Limites actuelles}

\begin{itemize}
  \item \textbf{Plan Finnhub gratuit} : seul le flux \emph{trade} est diffusé en
        temps réel ; quotes et candles temps réel nécessitent un abonnement.
  \item \textbf{Métadonnées symboles} : secteur et nom de société absents (non
        fournis par le plan gratuit).
  \item \textbf{Entrepôt} : DuckDB local remplace Snowflake pour rester gratuit ;
        le code SQL reste néanmoins portable.
  \item \textbf{Environnement de test} : saturation disque limitant les essais
        de bout en bout conteneurisés.
\end{itemize}

\subsection{Améliorations envisageables}

\begin{itemize}
  \item Schema Registry et format Avro pour le contrat des messages.
  \item Traitement de flux (Kafka Streams, Spark Streaming).
  \item Qualité de données automatisée (Great Expectations).
  \item Supervision (Prometheus, Grafana).
  \item Intégration et déploiement continus (CI/CD via GitHub Actions).
  \item Bascule vers un entrepôt cloud (Snowflake) en réutilisant les modèles dbt.
\end{itemize}

% ==================================================================
\section{Conclusion}
% ==================================================================

Le projet met en œuvre une pipeline de données \textbf{complète et fonctionnelle},
du streaming temps réel jusqu'à la modélisation analytique. Les cinq étapes ont
été réalisées et validées : ingestion Finnhub vers Kafka, stockage en data lake
structuré en trois couches, orchestration par Airflow, modélisation dbt, et
conteneurisation Docker.

Au-delà de la simple réalisation, la démarche a privilégié la \textbf{rigueur
d'ingénierie} : décisions de conception explicites et justifiées, validation par
exécution réelle à chaque étape, gestion de version structurée et documentation
systématique. L'architecture obtenue est modulaire et extensible, prête à
accueillir les améliorations envisagées (qualité de données, supervision,
entrepôt cloud).

\end{document}
